% Figure: the quasi-static regime map. System size L against frequency f, with
% the diagonal L = c/f marking the wavelength. Below it, lumped/quasi-static;
% above it, full-wave. A second line at L = lambda/10 marks the practical
% lumped boundary. Log-log axes.
\begin{figure}[htbp]
\centering
\begin{tikzpicture}
\begin{loglogaxis}[
  width=12cm, height=8cm,
  xlabel={frequency $f$ (Hz)},
  ylabel={system size $L$ (m)},
  xmin=1, xmax=1e12, ymin=1e-4, ymax=1e6,
  axis lines=left,
  tick label style={font=\scriptsize},
  label style={font=\small},
  legend pos=south west,
  legend style={font=\scriptsize},
]
% wavelength line L = c/f, c = 3e8
\addplot[engblue, very thick, samples=80, domain=1:1e12] {3e8/x};
\addlegendentry{$L=\lambda=c/f$}
% a tenth of a wavelength: practical lumped boundary
\addplot[engred, thick, dashed, samples=80, domain=1:1e12] {3e7/x};
\addlegendentry{$L=\lambda/10$}
% region labels
\node[enggray, font=\scriptsize, align=center] at (axis cs:1e3,1e-2)
  {quasi-static\\(lumped circuit)};
\node[engblack, font=\scriptsize, align=center] at (axis cs:1e9,1e3)
  {full-wave\\(fields propagate)};
% mark some devices
\filldraw[engamber] (axis cs:50,0.3) circle [radius=2pt];
\node[engamber, font=\scriptsize, anchor=west] at (axis cs:80,0.3)
  {mains transformer};
\filldraw[engamber] (axis cs:1e9,0.1) circle [radius=2pt];
\node[engamber, font=\scriptsize, anchor=west] at (axis cs:1.4e9,0.1)
  {GHz board trace};
\filldraw[engamber] (axis cs:1e8,30) circle [radius=2pt];
\node[engamber, font=\scriptsize, anchor=east] at (axis cs:7e7,30)
  {FM dipole};
\end{loglogaxis}
\end{tikzpicture}
\caption[Quasi-static regime map]{The boundary between the quasi-static
regime, where Maxwell's equations reduce to circuit laws, and the full-wave
regime, where the fields propagate as waves. The dividing line is the
wavelength $\lambda = c/f$ set against the system size $L$. A device much
smaller than a wavelength (below the blue line) is well described by lumped
capacitance, inductance, and resistance; a device comparable to or larger
than a wavelength must be treated by the full equations. The practical
working boundary sits near $L=\lambda/10$ (red dashed), where propagation
delay across the device reaches a tenth of a period. A mains transformer at
$50\,\si{\hertz}$ sits deep in the quasi-static corner; a gigahertz board
trace sits on the full-wave side of the same plot.}
\label{fig:vol04ch08:quasistatic-regimes}
\end{figure}
